\chapter{Yüklemler ve Altküme Operatörleri}
\label{ch:predicates_operators}

%% ============================================================
\section{Konu}
%% ============================================================

\subsection{Altküme Yüklemleri}

$(X, \tau)$ bir topolojik uzay ve $A \subseteq X$ olsun.

\begin{table}[h]
\centering
\begin{tabular}{ll}
\toprule
\textbf{Yüklem} & \textbf{Tanım} \\
\midrule
Açık (open) & $A \in \tau$ \\
Kapalı (closed) & $X \setminus A \in \tau$ \\
Clopen & $A \in \tau$ ve $X \setminus A \in \tau$ \\
Yoğun (dense) & Her boş olmayan $U \in \tau$ için $U \cap A \neq \emptyset$ \\
Hiçbiryerde-yoğun & $\mathrm{int}(\overline{A}) = \emptyset$ \\
\bottomrule
\end{tabular}
\caption{Altküme yüklemleri}
\end{table}

\subsection{Altküme Operatörleri}

\begin{table}[h]
\centering
\begin{tabular}{lll}
\toprule
\textbf{Operatör} & \textbf{Gösterim} & \textbf{Tanım} \\
\midrule
Kapanış & $\overline{A} = \mathrm{cl}(A)$ & $A$'yı içeren en küçük kapalı küme \\
İç & $A^\circ = \mathrm{int}(A)$ & $A$'ya dahil en büyük açık küme \\
Sınır & $\partial A = \mathrm{bd}(A)$ & $\overline{A} \setminus A^\circ$ \\
Türev kümesi & $A' = d(A)$ & $A$'nın birikme noktaları kümesi \\
\bottomrule
\end{tabular}
\caption{Altküme operatörleri}
\end{table}

\noindent\textbf{Birikme noktası:} $x \in A'$ $\iff$ her $U \ni x$ açık için
$U \cap (A \setminus \{x\}) \neq \emptyset$.

\begin{sezgi}
Bir $A$ kümesini, $X$ üzerinde duran bir ``boya lekesi'' gibi düşünün.
\textbf{İç} $A^\circ$, lekenin tamamen içine sığabildiğiniz açık bir
komşuluğunuzun olduğu noktalardır --- ``rahatça içerideyim'' bölgesi.
\textbf{Kapanış} $\mathrm{cl}(A)$, lekeye dokunan her noktayı ekler:
leke artı onun ``gölgesi''. \textbf{Sınır} $\partial A$ ise tam kenardır:
her açık komşuluğu hem $A$'dan hem $X\setminus A$'dan kesen noktalar.
Bu üçlü, $X$'i ayrık üç parçaya böler: iç, sınır ve \textbf{dış}
$\mathrm{ext}(A) = (X\setminus A)^\circ$.
\end{sezgi}

\begin{figure}[ht]
  \centering
  \begin{tikzpicture}[font=\small,
    nokta/.style={circle, draw, minimum size=7mm, inner sep=1pt},
    ic/.style={nokta, fill=green!18},
    sin/.style={nokta, fill=orange!22},
    dis/.style={nokta, fill=blue!10}]
  % X = {1,2,3,4,5}; A = {1,2,3}
  % int(A)={1,2} (yesil), bd(A)={3} (turuncu), ext(A)={4,5} (mavi)
  \node[ic]  (n1) at (0,0)   {$1$};
  \node[ic]  (n2) at (1.5,0) {$2$};
  \node[sin] (n3) at (3,0)   {$3$};
  \node[dis] (n4) at (4.5,0) {$4$};
  \node[dis] (n5) at (6,0)   {$5$};

  % A kapanis cercevesi cl(A) = {1,2,3}
  \draw[red!70!black, thick, rounded corners]
    ($(n1.north west)+(-0.25,0.55)$) rectangle ($(n3.south east)+(0.25,-0.55)$);
  \node[red!70!black] at (1.5,1.05) {$\mathrm{cl}(A)=\{1,2,3\}$};

  % int(A) cercevesi
  \draw[green!55!black, thick, dashed, rounded corners]
    ($(n1.north west)+(-0.15,0.3)$) rectangle ($(n2.south east)+(0.15,-0.3)$);

  % Etiketler
  \node[green!45!black, font=\scriptsize] at (0.75,-1.0) {iç: $A^\circ=\{1,2\}$};
  \node[orange!75!black, font=\scriptsize] at (3,-1.0) {sınır: $\partial A=\{3\}$};
  \node[blue!65!black, font=\scriptsize] at (5.25,-1.0) {dış: $\mathrm{ext}(A)=\{4,5\}$};

  \node[font=\scriptsize, align=center] at (3,-1.7)
    {$X = A^\circ \;\cup\; \partial A \;\cup\; \mathrm{ext}(A)$ (ayrık parçalanış)};
\end{tikzpicture}

  \caption{$A=\{1,2,3\}$ için iç $A^\circ=\{1,2\}$ (yeşil), sınır
    $\partial A=\{3\}$ (turuncu) ve dış $\mathrm{ext}(A)=\{4,5\}$ (mavi);
    kapanış $\mathrm{cl}(A)$ kırmızı çerçevedir. $X$ bu üç parçaya ayrık olarak bölünür.}
  \label{fig:ch05:ic_kapanis_sinir}
\end{figure}

\subsection{Komşuluk Sistemi}

\begin{definition}[Komşuluk Sistemi]
$x \in X$ için:
\[
  N(x) = \{N \subseteq X : \exists U \in \tau,\, x \in U \subseteq N\}
\]
\end{definition}

\noindent\textbf{Komşuluk Aksiyomları (N1--N4):}
\begin{enumerate}
  \item[(N1)] $x \in N$, her $N \in N(x)$ için
  \item[(N2)] $X \in N(x)$
  \item[(N3)] $N_1, N_2 \in N(x) \Rightarrow N_1 \cap N_2 \in N(x)$
  \item[(N4)] $N \in N(x)$, $N \subseteq M \Rightarrow M \in N(x)$
\end{enumerate}

%% ============================================================
\section{Teoremler}
%% ============================================================

\begin{theorem}[Kuratowski Kapanış Aksiyomları]
\label{thm:kuratowski}
$\mathrm{cl}: \mathcal{P}(X) \to \mathcal{P}(X)$ operatörü şu dört özelliği sağlar:
\begin{enumerate}
  \item[(K1)] $\mathrm{cl}(\emptyset) = \emptyset$
  \item[(K2)] $A \subseteq \mathrm{cl}(A)$
  \item[(K3)] $\mathrm{cl}(\mathrm{cl}(A)) = \mathrm{cl}(A)$
  \item[(K4)] $\mathrm{cl}(A \cup B) = \mathrm{cl}(A) \cup \mathrm{cl}(B)$
\end{enumerate}
Bu dört özelliği sağlayan her operatör bir topoloji tanımlar.
\end{theorem}

\begin{figure}[ht]
  \centering
  \begin{tikzpicture}[font=\small, node distance=7mm,
    aks/.style={draw, rounded corners, fill=blue!6, inner sep=6pt,
                text width=4.6cm, align=center}]
  \node[aks] (k1) {\textbf{(K1)}\\ $\mathrm{cl}(\emptyset)=\emptyset$};
  \node[aks, right=of k1] (k2) {\textbf{(K2)}\\ $A \subseteq \mathrm{cl}(A)$};
  \node[aks, below=of k1] (k3) {\textbf{(K3)}\\ $\mathrm{cl}(\mathrm{cl}(A))=\mathrm{cl}(A)$};
  \node[aks, below=of k2] (k4) {\textbf{(K4)}\\ $\mathrm{cl}(A\cup B)=\mathrm{cl}(A)\cup \mathrm{cl}(B)$};

  \node[draw, rounded corners, fill=green!10, inner sep=6pt,
        below=10mm of $(k3)!0.5!(k4)$, text width=7.2cm, align=center] (sonuc)
    {Bu dört aksiyomu sağlayan her $\mathrm{cl}$ operatörü\\
     \emph{tek bir} topoloji belirler: $\tau=\{X\setminus \mathrm{cl}(A)\}$.};

  \draw[-{Stealth[length=2.4mm]}, thick] (k3.south) -- (k3.south |- sonuc.north);
  \draw[-{Stealth[length=2.4mm]}, thick] (k4.south) -- (k4.south |- sonuc.north);
\end{tikzpicture}

  \caption{Kuratowski aksiyomları K1--K4 ve sonucu: bu dördünü sağlayan her
    $\mathrm{cl}$ operatörü tek bir topoloji belirler.}
  \label{fig:ch05:kuratowski}
\end{figure}

\begin{proof}[İspat eskizi]
Kapanış, $A$'yı içeren tüm kapalı kümelerin kesişimidir:
$\mathrm{cl}(A) = \bigcap\{C : C \text{ kapalı},\, A\subseteq C\}$.
\textbf{(K1)} $\emptyset$ zaten kapalı, en küçük kapalı üst-küme kendisidir.
\textbf{(K2)} her $C\supseteq A$ içerdiğinden $A\subseteq \mathrm{cl}(A)$.
\textbf{(K3)} $\mathrm{cl}(A)$ kapalıdır; kapalı bir kümenin kapanışı kendisidir
(idempotentlik). \textbf{(K4)} $\subseteq$ yönü monotonluktan, $\supseteq$ yönü
$\mathrm{cl}(A)\cup\mathrm{cl}(B)$'nin $A\cup B$'yi içeren kapalı bir küme olmasından gelir.
Sonlu uzaylarda dördü de \texttt{kuratowski\_closure\_check} ile doğrulanır (Örnek~2).
\end{proof}

\begin{theorem}[İç-Kapanış Dualitesi]
\label{thm:interior_closure_duality}
Her $A \subseteq X$ için:
\[
  A^\circ = X \setminus \mathrm{cl}(X \setminus A), \qquad
  \mathrm{cl}(A) = X \setminus \mathrm{int}(X \setminus A).
\]
\end{theorem}

\begin{proof}[İspat eskizi]
$A^\circ$ tanım gereği $A$'ya dahil \emph{en büyük açık} kümedir; $\mathrm{cl}(B)$ ise
$B$'yi içeren \emph{en küçük kapalı} kümedir. Tümleme alma açık~$\leftrightarrow$~kapalı
ve ``en büyük''~$\leftrightarrow$~``en küçük'' ile $\subseteq$ ilişkisini ters çevirir.
$B=X\setminus A$ alın: $X\setminus A$'yı içeren en küçük kapalı kümenin tümleyeni,
$A$'ya dahil en büyük açık kümedir; yani $A^\circ = X\setminus\mathrm{cl}(X\setminus A)$.
İkinci eşitlik $A\mapsto X\setminus A$ ikamesiyle simetrik olarak çıkar.
\end{proof}

\begin{theorem}[Komşuluk Sistemi Round-Trip]
\label{thm:neighborhood_round_trip}
N1--N4'ü sağlayan her $\{N(x)\}_{x \in X}$ ailesi,
$\tau = \{U : \forall x \in U, U \in N(x)\}$ şeklinde tek bir topolojiyi geri üretir.
\end{theorem}

\begin{proof}[İspat eskizi]
İleri yön: $\tau$'dan $N(x)=\{N : \exists U\in\tau,\, x\in U\subseteq N\}$ tanımlayın.
(N1) $x\in U\subseteq N$ olduğundan $x\in N$. (N2) $X\in\tau$ her noktayı içerir.
(N3) $U_1\cap U_2\in\tau$ açık olduğundan kesişim de komşuluktur. (N4) üst-küme almak
tanımı bozmaz. Geri yön: $\tau=\{U : \forall x\in U,\, U\in N(x)\}$ ailesinin topoloji
aksiyomlarını sağladığı ve başlangıç sistemini geri verdiği doğrulanır; teklik bu
eşleşmenin tersinir olmasından gelir.
\end{proof}

%% ============================================================
\section{Algoritmalar}
%% ============================================================

\begin{algorithm}
\caption{Sonlu Kapanış Hesabı}
\begin{algorithmic}[1]
\Require $A \subseteq X$, $\tau$
\Ensure $\mathrm{cl}(A)$
\State $\textit{closed} \leftarrow \{X \setminus U : U \in \tau\}$
\State $\textit{result} \leftarrow X$
\For{each $C \in \textit{closed}$}
    \If{$A \subseteq C$ \textbf{and} $|C| < |\textit{result}|$}
        \State $\textit{result} \leftarrow C$
    \EndIf
\EndFor
\Return $\textit{result}$
\end{algorithmic}
\end{algorithm}

\noindent\textbf{Karmaşıklık:} $O(|\tau| \cdot |X|)$.

\begin{algorithm}
\caption{Türev Kümesi Hesabı}
\begin{algorithmic}[1]
\Require $A \subseteq X$, $\tau$
\Ensure $d(A)$
\State $\textit{acc} \leftarrow \emptyset$
\For{each $x \in X$}
    \If{her açık $U \ni x$ için $U \cap (A \setminus \{x\}) \neq \emptyset$}
        \State $\textit{acc}.\text{add}(x)$
    \EndIf
\EndFor
\Return $\textit{acc}$
\end{algorithmic}
\end{algorithm}

\noindent\textbf{Karmaşıklık:} $O(|X| \cdot |\tau|)$.

%% ============================================================
\section{pytop API}
%% ============================================================

\begin{lstlisting}[language=Python, caption={Subset operatörleri importları}]
from pytop import (
    analyze_predicate,
    closure_of_subset, interior_of_subset, boundary_of_subset,
    derived_set_of_subset,
    is_open_subset, is_closed_subset, is_dense_subset, is_nowhere_dense_subset,
    neighborhood_system, character_at_point, analyze_neighborhood_system,
)
\end{lstlisting}

\begin{table}[h]
\centering
\begin{tabular}{lll}
\toprule
\textbf{Fonksiyon} & \textbf{İmza} & \textbf{Açıklama} \\
\midrule
\texttt{analyze\_predicate} & \texttt{(space, predicate, subset)} & Yüklem sorgula \\
\texttt{closure\_of\_subset} & \texttt{(space, subset)} & Kapanış \\
\texttt{interior\_of\_subset} & \texttt{(space, subset)} & İç \\
\texttt{boundary\_of\_subset} & \texttt{(space, subset)} & Sınır \\
\texttt{derived\_set\_of\_subset} & \texttt{(space, subset)} & Türev kümesi \\
\texttt{neighborhood\_system} & \texttt{(carrier, topology, point)} & $N(x)$ \\
\texttt{character\_at\_point} & \texttt{(carrier, topology, point)} & $\chi(x)$ \\
\texttt{analyze\_neighborhood\_system} & \texttt{(carrier, topology, point=None)} & Tam analiz \\
\bottomrule
\end{tabular}
\caption{Altküme operatör fonksiyonları}
\end{table}

\noindent\textbf{Not:} \texttt{neighborhood\_system} ve ilgili fonksiyonlar \texttt{carrier} (liste)
ve \texttt{topology} (liste) alır. Space nesnesi için: \texttt{list(space.carrier)},
\texttt{list(space.topology)}.

%% ============================================================
\section{Örnekler}
%% ============================================================

\subsection{Örnek 1: Kapanış ve İç — Sierpiński Uzayı}

\begin{lstlisting}[language=Python]
from pytop import sierpinski_space, closure_of_subset, interior_of_subset, boundary_of_subset

s = sierpinski_space()
print("cl({0})  =", closure_of_subset(s, {0}).value)
print("cl({1})  =", closure_of_subset(s, {1}).value)
print("int({0}) =", interior_of_subset(s, {0}).value)
print("int({1}) =", interior_of_subset(s, {1}).value)
print("bd({1})  =", boundary_of_subset(s, {1}).value)
\end{lstlisting}

\begin{lstlisting}[caption={Çıktı}]
cl({0})  = frozenset({0})
cl({1})  = frozenset({0, 1})
int({0}) = frozenset()
int({1}) = frozenset({1})
bd({1})  = frozenset({0})
\end{lstlisting}

$\{0\}$'ın kapanışı kendisidir: $\{0\} = X \setminus \{1\}$ zaten kapalıdır.
$\{1\}$'in kapanışı $X$'tir: $\{1\}$ kapalı değildir.

\subsection{Örnek 2: Türev Kümesi}

\begin{lstlisting}[language=Python]
from pytop import sierpinski_space, derived_set_of_subset, is_nowhere_dense_subset

s = sierpinski_space()
print("d({0}) =", derived_set_of_subset(s, {0}).value)
print("d({1}) =", derived_set_of_subset(s, {1}).value)
print("{0} nowhere dense? =", is_nowhere_dense_subset(s, {0}).status)
\end{lstlisting}

\begin{lstlisting}[caption={Çıktı}]
d({0}) = frozenset()
d({1}) = frozenset({0})
{0} nowhere dense? = true
\end{lstlisting}

$d(\{1\}) = \{0\}$: $0$'ı içeren her $\{0,1\}$, $\{1\} \setminus \{0\} = \{1\}$ ile kesişir.

\subsection{Örnek 3: Komşuluk Sistemi}

\begin{lstlisting}[language=Python]
from pytop import sierpinski_space, neighborhood_system, character_at_point

s = sierpinski_space()
carrier = list(s.carrier)
topology = list(s.topology)
print("N(0) =", neighborhood_system(carrier, topology, 0).value)
print("N(1) =", neighborhood_system(carrier, topology, 1).value)
print("chi(0) =", character_at_point(carrier, topology, 0).value)
print("chi(1) =", character_at_point(carrier, topology, 1).value)
\end{lstlisting}

\begin{lstlisting}[caption={Çıktı}]
N(0) = [['0', '1']]
N(1) = [['1'], ['0', '1']]
chi(0) = 1
chi(1) = 2
\end{lstlisting}

$N(0) = \{X\}$: $0$ yalnızca $X$'te açık. $N(1) = \{\{1\}, X\}$: yerel baz büyüklükleri
$\chi(0)=1$, $\chi(1)=2$.

\subsection{Örnek 4: İç / Kapanış / Sınır / Dış Parçalanışı}

\begin{lstlisting}[language=Python]
from pytop import (make_topology, interior_of_subset, closure_of_subset,
                   boundary_of_subset, exterior_of_subset)

sp = make_topology({1, 2, 3, 4, 5}, {1, 2}, {4, 5}, {1, 2, 4, 5})
A = {1, 2, 3}
print("int(A) =", interior_of_subset(sp, A).value)
print("cl(A)  =", closure_of_subset(sp, A).value)
print("bd(A)  =", boundary_of_subset(sp, A).value)
print("ext(A) =", exterior_of_subset(sp, A).value)
\end{lstlisting}

\begin{lstlisting}[caption={Çıktı}]
int(A) = frozenset({1, 2})
cl(A)  = frozenset({1, 2, 3})
bd(A)  = frozenset({3})
ext(A) = frozenset({4, 5})
\end{lstlisting}

İç $\{1,2\}$, sınır $\{3\}$ ve dış $\{4,5\}$ birlikte $X$'i ayrık olarak kaplar:
$X = A^\circ \cup \partial A \cup \mathrm{ext}(A)$ (Şekil~\ref{fig:ch05:ic_kapanis_sinir}).

\subsection{Örnek 5: Kuratowski Aksiyomlarının Doğrulanması}

\begin{lstlisting}[language=Python]
from pytop import finite_chain_space, kuratowski_closure_check

s = finite_chain_space(3)
report = kuratowski_closure_check(list(s.carrier), list(s.topology))
for k in sorted(report):
    print(f"{k:<13}= {report[k]}")
\end{lstlisting}

\begin{lstlisting}[caption={Çıktı}]
all          = True
empty        = True
extensive    = True
finite_union = True
idempotent   = True
\end{lstlisting}

\texttt{empty}~$\to$~(K1), \texttt{extensive}~$\to$~(K2), \texttt{idempotent}~$\to$~(K3),
\texttt{finite\_union}~$\to$~(K4); \texttt{all} hepsinin sağlandığını özetler
(Şekil~\ref{fig:ch05:kuratowski}).

\subsection{Örnek 6: Yoğun Nokta ve Birikme Kümesi}

\begin{lstlisting}[language=Python]
from pytop import finite_chain_space, closure_of_subset, is_dense_subset, derived_set_of_subset

c = finite_chain_space(4)
print("cl({1})    =", closure_of_subset(c, {1}).value)
print("{1} dense? =", is_dense_subset(c, {1}).status)
print("d({3})     =", derived_set_of_subset(c, {3}).value)
print("cl({4})    =", closure_of_subset(c, {4}).value)
\end{lstlisting}

\begin{lstlisting}[caption={Çıktı}]
cl({1})    = frozenset({1, 2, 3, 4})
{1} dense? = true
d({3})     = frozenset({4})
cl({4})    = frozenset({4})
\end{lstlisting}

$1$ jeneriktir ($\mathrm{cl}(\{1\})=X$, yoğun), $4$ kapalı bir noktadır
($\mathrm{cl}(\{4\})=\{4\}$); Şekil~\ref{fig:ch05:yogunluk}.

\begin{figure}[ht]
  \centering
  \begin{tikzpicture}[font=\small,
    nokta/.style={circle, draw, minimum size=7mm, inner sep=1pt, fill=gray!8},
    yogun/.style={circle, draw, minimum size=7mm, inner sep=1pt, fill=green!22}]
  % finite_chain_space(4): acik kumeler {1} subset {1,2} subset {1,2,3} subset X
  % cl({1}) = X  -> 1 yogun nokta;  cl({4}) = {4} -> 4 kapali nokta
  \node[yogun] (n1) at (0,0)   {$1$};
  \node[nokta] (n2) at (2,0)   {$2$};
  \node[nokta] (n3) at (4,0)   {$3$};
  \node[nokta] (n4) at (6,0)   {$4$};

  \foreach \a/\b in {n1/n2, n2/n3, n3/n4}
    \draw[-{Stealth[length=2.4mm]}, thick] (\a) -- (\b);

  \node[font=\scriptsize] at (1,0.45) {özelleşme};

  % cl({1}) = X kapanis cercevesi
  \draw[green!55!black, thick, rounded corners]
    ($(n1.north west)+(-0.25,0.95)$) rectangle ($(n4.south east)+(0.25,-0.3)$);
  \node[green!45!black, font=\scriptsize] at (3,1.25) {$\mathrm{cl}(\{1\})=X$ — $\{1\}$ yoğun};

  \node[font=\scriptsize, align=center] at (6,-0.95)
    {$\mathrm{cl}(\{4\})=\{4\}$\\ kapalı nokta};
  \node[font=\scriptsize, align=center] at (0,-0.95)
    {jenerik\\ (yoğun)};
\end{tikzpicture}

  \caption{\texttt{finite\_chain\_space(4)} zincirinde özelleşme oku: $1$ jenerik
    (yoğun) noktadır, $\mathrm{cl}(\{1\})=X$; $4$ ise kapalı noktadır, $\mathrm{cl}(\{4\})=\{4\}$.}
  \label{fig:ch05:yogunluk}
\end{figure}

\begin{karsiornek}
``İç küme her zaman kapalıdır'' yanılgısı. Örnek~4'te $A^\circ=\{1,2\}$ açıktır ama
\textbf{kapalı değildir} ($X\setminus\{1,2\}=\{3,4,5\}$ açık değil). İç ve dış daima
açıktır; kapanış ve sınır daima kapalıdır --- bu iki aileyi karıştırmak en sık yapılan
hatadır. Ayrıca $\partial A=\emptyset$ ancak ve ancak $A$ clopen ise olur: Örnek~``Sınır
Hesabı''nda açık $\{1\}$ için $\partial\{1\}=\emptyset$, fakat clopen olmayan $\{1,2\}$
için sınır boş değildir.
\end{karsiornek}

%% ============================================================
\section{Alıştırmalar}
%% ============================================================

\subsection*{Kodlama Alıştırmaları}
\begin{enumerate}[label=\textbf{K\arabic*.}]
  \item İndirgenmiş iki-noktalı topolojide her noktanın komşuluk sistemini hesaplayın.
  \item $X = \{1,2,3,4\}$, $\tau = \{\emptyset, \{1,2\}, \{3,4\}, X\}$ topolojisinde
        $\mathrm{bd}(\{1,2,3\})$ ve $\mathrm{cl}(\{1,2,3\})$ hesaplayın.
  \item \texttt{finite\_chain\_space(4)} zincirinde $\{1,2\}$'nin iç, kapanış ve sınır
        kümelerini bulun.
  \item \texttt{make\_topology(\{1,2,3,4,5\}, \{1,2\}, \{4,5\}, \{1,2,4,5\})} uzayında
        $A=\{1,2,3\}$ için \texttt{interior\_of\_subset}, \texttt{boundary\_of\_subset} ve
        \texttt{exterior\_of\_subset} çalıştırıp $X = A^\circ \cup \partial A \cup
        \mathrm{ext}(A)$ ayrık parçalanışını doğrulayın.
  \item \texttt{finite\_chain\_space(4)} için \texttt{kuratowski\_closure\_check} çağırıp
        \texttt{all} alanının \texttt{True} döndüğünü gösterin; ardından $\{1\}$'in yoğun
        olduğunu \texttt{is\_dense\_subset} ile teyit edin.
\end{enumerate}

\subsection*{Teori Alıştırmaları}
\begin{enumerate}[label=\textbf{T\arabic*.}]
  \item $A^\circ = X \setminus \mathrm{cl}(X \setminus A)$ eşitliğini Kuratowski
        aksiyomlarını kullanarak kanıtlayın.
  \item $A$'nın yoğun olması için gerek ve yeter koşulun $\mathrm{cl}(A) = X$ olduğunu
        gösterin.
  \item Her $A \subseteq X$ için $X = A^\circ \cup \partial A \cup \mathrm{ext}(A)$
        olduğunu ve üç kümenin ikişer ikişer ayrık olduğunu ispatlayın.
        \ipucu{$\mathrm{ext}(A)=(X\setminus A)^\circ$ ve $\partial A=\mathrm{cl}(A)\setminus A^\circ$.}
  \item $A^\circ$ kümesinin daima açık olduğunu, fakat genel olarak kapalı
        \textbf{olmadığını} gösteren bir karşı-örnek verin. \ipucu{Örnek~4'teki $\{1,2\}$.}
\end{enumerate}
